\documentclass[10pt,twocolumn]{article}
\usepackage[T1]{fontenc}
\usepackage[utf8]{inputenc}
\usepackage{lmodern}
\usepackage[margin=1.8cm,top=2.4cm,bottom=2cm,columnsep=0.6cm]{geometry}
\usepackage{fancyhdr}
\usepackage{titlesec}
\usepackage{booktabs}
\usepackage{amsmath,amssymb,amsfonts}
\usepackage{microtype}
\usepackage{xcolor}
\usepackage{mdframed}
\usepackage{parskip}
\usepackage{enumitem}
\usepackage{hyperref}

\definecolor{rulered}{RGB}{180,30,30}
\definecolor{boxgray}{RGB}{245,245,245}
\definecolor{keywordgray}{RGB}{100,100,100}

\pagestyle{fancy}
\fancyhf{}
\fancyhead[L]{\textsc{\small Claude Mag}}
\fancyhead[C]{\small\textit{Machine Learning Research Digest}}
\fancyhead[R]{\small 2026-07-22}
\fancyfoot[C]{\small\thepage}
\renewcommand{\headrulewidth}{0.8pt}

\titleformat{\section}{\normalsize\bfseries\color{rulered}}{}{0pt}{}%
  [\vspace{-4pt}\color{rulered}\rule{\columnwidth}{0.4pt}]
\titlespacing{\section}{0pt}{8pt}{4pt}

\hypersetup{colorlinks=true,urlcolor=rulered,linkcolor=black}

\begin{document}
\setlength{\parindent}{0pt}
\setlength{\parskip}{4pt}

{\small\color{keywordgray}\textbf{Keywords:}
  pretraining \textbullet{}
  reinforcement learning \textbullet{}
  scaling laws \textbullet{}
  chess reasoning \textbullet{}
  post-training}\par\vspace{4pt}

{\LARGE\bfseries\raggedright
  Understanding Reasoning from Pretraining to Post-Training}\par\vspace{4pt}

{\large\itshape\raggedright
  NYU and Collaborators Use Chess as a Controlled Testbed to Uncover a
  Joint Pretraining-to-RL Scaling Law and Show That RL Concentrates
  Rather Than Creates Competence}\par\vspace{6pt}

{\small\textbf{By} Jingyan Shen, Ang Li, Salman Rahman, Yifan Sun,
  Micah Goldblum, Matus Telgarsky, Pavel Izmailov
  (New York University, Modal Labs, UCLA, UIUC, Columbia University)}\par\vspace{2pt}

{\small\color{keywordgray}arXiv:2607.16097 \textbullet{}
  July 17, 2026}
\par\vspace{2pt}\noindent{\color{rulered}\rule{\columnwidth}{0.6pt}}\vspace{6pt}

Reinforcement learning post-training has become a core ingredient in
every state-of-the-art language model, yet the field lacks basic
answers to two questions: how do pretraining choices shape the returns
to RL compute, and what does RL actually do to the model?
A team from NYU, Modal Labs, UCLA, UIUC, and Columbia now answers
both questions using chess --- a fully observable reasoning domain where
every move can be evaluated with ground-truth certainty.

\begin{mdframed}[backgroundcolor=boxgray,linecolor=rulered,linewidth=1pt,
    innerleftmargin=6pt,innerrightmargin=6pt,
    innertopmargin=6pt,innerbottommargin=6pt]
\textbf{\small AT A GLANCE}\par\vspace{4pt}
\begin{description}[leftmargin=0pt,labelsep=4pt,itemsep=2pt,parsep=0pt,
    font=\small\bfseries]
  \item[Problem:] RL post-training is studied in isolation from pretraining,
    leaving the joint compute allocation problem and RL's mechanistic
    role undefined.
  \item[Why it matters:] Understanding the pretraining-to-RL interface
    is necessary for principled scaling decisions and for knowing what
    RL can and cannot add to a model.
  \item[Method:] Pretraining of 5M--1B parameter models on chess games,
    SFT on synthetic reasoning traces, RL on chess puzzles; move-level
    probability analysis across checkpoints.
  \item[Result:] A joint scaling law $R(L_{\text{pre}}, C_{\text{RL}})$
    fits held-out runs accurately; RL concentrates probability mass on
    high-value moves already present in the pretrained distribution.
  \item[Evidence:] Scaling law extrapolates accurately to a 1B-parameter
    model with 10$\times$ more RL compute than any training run used
    to fit it.
\end{description}
\end{mdframed}

\section{The Big Picture}

Reinforcement learning post-training is now applied to every frontier
model, yet the field studies it in isolation: RL is treated as a
plug-in module applied to an already-trained base model, with no
principled account of how the two training stages interact.
This isolation has produced two blind spots.
First, practitioners cannot predict how much RL compute a given
pretrained model will reward --- leading to costly trial-and-error
compute allocation.
Second, there is no controlled mechanistic account of what RL
actually modifies in the model's representations and output
distributions; the prevalent intuition is that ``RL teaches
reasoning'' but this is rarely examined at the level of
individual predictions.

The paper uses chess as a controlled testbed because it offers ground
truth for every decision, a known policy complexity profile, and
a dataset of human games spanning all skill levels ---
ideal conditions for clean scaling experiments that would be
impossible on open-ended language tasks.

\section{Why Existing Approaches Fall Short}

Prior work on RL for reasoning studies either the RL stage in
isolation (holding pretraining fixed) or varies pretraining at large
scale without controlling the RL stage.
Neither approach can identify the joint interaction between the two
stages.

Studying RL in isolation conflates the contribution of the pretrained
representation with the contribution of RL optimization, making it
impossible to determine whether a performance gain came from RL
itself or from the quality of the starting point.
Conversely, varying pretraining at fixed RL compute answers which base
model RL benefits most, but cannot characterize the functional form
of that benefit or reveal what RL does to the model's
internal decision-making.

\section{Core Mechanism}

The authors vary both pretraining quality (model size and data volume)
and RL compute in a factorial design, measuring post-RL reward for
each combination.
They fit a parametric scaling law to the resulting data and validate
it on held-out (pretraining quality, RL compute) pairs.
In parallel, they analyze the move-probability distributions of each
checkpoint to track how RL reshapes the policy beyond aggregate reward.

\section{Technical Formulation}

The joint scaling law expresses post-RL reward $R$ as a function of
pretraining loss $L_{\text{pre}}$ and RL compute $C_{\text{RL}}$:
\[
R(L_{\text{pre}},\, C_{\text{RL}})
  = R_\infty - \phi(L_{\text{pre}}) \cdot C_{\text{RL}}^{-\gamma}
\]
where $R_\infty$ is the asymptotic reward ceiling, $\phi(L_{\text{pre}})$
is a monotonically increasing function of pretraining loss
(worse pretraining shifts the ceiling down and slows RL convergence),
and $\gamma$ is a universal power-law exponent estimated from data.

The function $\phi$ is fit empirically and found to be approximately
linear in $L_{\text{pre}}$ over the studied range, giving:
\[
\phi(L_{\text{pre}}) \approx \alpha \cdot L_{\text{pre}} + \beta
\]
with $\alpha$ and $\beta$ estimated jointly with $\gamma$ via nonlinear
least squares.
This form implies that the marginal value of RL compute is linear in
pretraining loss: every unit improvement in pretraining quality
proportionally increases the RL compute budget required to reach
any given reward threshold.

\section{Learning Procedure}

Training proceeds in three stages applied to each factorial combination:

\begin{enumerate}[itemsep=2pt,parsep=0pt]
  \item \textbf{Pretraining:} Next-token prediction on Lichess PGN
    game records.
    Models range from 5M to 1B parameters; data volume is varied
    orthogonally from 10M to 2B tokens.
  \item \textbf{Supervised fine-tuning:} Training on synthetically
    generated chain-of-thought traces that articulate candidate
    move rationale in natural language.
  \item \textbf{RL post-training:} GRPO with binary correctness reward
    on chess puzzles; no intermediate-step reward.
    RL compute is varied from $10^{18}$ to $10^{20}$ FLOPs across
    factorial conditions.
\end{enumerate}

\section{Experimental Findings}

The joint scaling law fits held-out (pretraining, RL compute) pairs
with relative error below 3\% across all studied configurations,
confirming that the parametric form captures the essential
interaction between the two training stages.

\begin{center}
\small
\begin{tabular}{lcc}
\toprule
Condition & Predicted $R$ & Observed $R$ \\
\midrule
1B params, low RL & 0.81 & 0.80 \\
1B params, high RL & 0.89 & 0.91 \\
70M params, high RL & 0.74 & 0.73 \\
5M params, high RL & 0.51 & 0.52 \\
\bottomrule
\end{tabular}
\end{center}

Mechanistic analysis reveals that RL concentrates the probability
distribution: the probability of the ground-truth top move increases
by an average of 31\% after RL for positions where the pretrained
model already ranks it top-1, while the probability assigned to
moves not in the pretrained model's top-10 changes by less
than 0.3\% on average.

\section{Ablations and Interpretation}

\textbf{SFT vs.\ RL:} SFT on reasoning traces reaches a lower
performance ceiling than RL, suggesting that imitation of human
reasoning traces is insufficient and RL's self-play signal is
necessary to reach high puzzle accuracy.

\textbf{RL without SFT:} Applying RL directly to the pretrained
model (without SFT) converges more slowly and to a lower final
reward, confirming that the chain-of-thought format learned during
SFT provides a useful scaffold for RL exploration.

\textbf{Mechanistic result:} RL does not inject knowledge of moves
the pretrained model assigned near-zero probability to.
This has a direct implication: RL post-training cannot recover from
severe pretraining data gaps.
A model that has never seen a particular chess strategy in pretraining
will not learn it through RL, regardless of RL compute budget.

\textbf{Scaling extrapolation:} The scaling law predicts post-RL
reward for a 1B-parameter model trained with 10$\times$ more RL
compute than any run used to fit the law, with only 1.8\% relative
error, suggesting the law is reliable for planning compute allocation
at larger scales.

\section*{Reference}

Jingyan Shen, Ang Li, Salman Rahman, Yifan Sun, Micah Goldblum,
Matus Telgarsky, Pavel Izmailov.\
\textbf{Understanding Reasoning from Pretraining to Post-Training}.
arXiv:2607.16097, July 2026.
\url{https://arxiv.org/abs/2607.16097}

\end{document}
